
\section{Related Work}
% \subsection{Guidance Strategies in Diffusion Models}
% Various methods aim to improve guidance in diffusion models. \textbf{CFG++}~\citep{chung2024cfg++} treats guidance as a manifold-constrained inverse problem, while \textbf{CFG Schedulers}~\citep{xianalysis} and \textbf{Apply Guidance in Interval}~\citep{kynkaanniemi2024applying} optimize time-dependent strength. \textbf{ReCFG}~\citep{xia2025rectified} and \textbf{TFG}~\citep{ye2024tfg} focus on correcting expectation shifts. Recent approaches explore dynamic adaptation: $\mathbf{S^2}$\textbf{-Guidance}~\citep{chen2025s} uses stochastic block-dropping, \textbf{REG}~\citep{gao2025reg} optimizes scaled joint distributions, \textbf{FBG}~\citep{koulischer2025feedback} uses feedback on conditional informativeness, and \textbf{Foresight Guidance}~\citep{wang2025foresight} frames CFG as a fixed-point iteration with short inner loops. \textbf{Density Guidance}~\citep{zhao2025devil} introduces detail-aware generation by controlling the likelihood-detail trade-off. In Flow Matching models~\citep{esser2024scalingrectifiedflowtransformers, lipman2022flow, fan2025vchitect, liu2022flow, gao2024lumina}, \textbf{CFG-Zero*}~\citep{fan2025cfg0} compensates for velocity errors.  

\subsection{Guidance Strategies in Diffusion Models}

Various methods aim to improve guidance in diffusion models. \textbf{CFG++}~\citep{chungcfg++} treats guidance as a manifold-constrained inverse problem~\citep{karczewski2025spacetime}, while \textbf{CFG Schedulers}~\citep{xianalysis} and \textbf{Apply Guidance in Interval}~\citep{kynkaanniemi2024applying} optimize time-dependent strength; similarly, \textbf{Stage-Wise Dynamics}~\citep{jin2025stage} investigates varying guidance requirements across different generation stages. \textbf{ReCFG}~\citep{xia2025rectified} and \textbf{TFG}~\citep{ye2024tfg} focus on correcting expectation shifts, a goal shared by \textbf{Rectified CFG++}~\citep{saini2025rectified} and \textbf{CFG-EC}~\citep{yang2025cfg} which propose mechanisms to rectify guidance errors and improve consistency. Recent approaches explore dynamic adaptation: $\mathbf{S^2}$\textbf{-Guidance}~\citep{chen2025s} uses stochastic block-dropping, \textbf{REG}~\citep{gao2025reg} optimizes scaled joint distributions, \textbf{FBG}~\citep{koulischer2025feedback} uses feedback on conditional informativeness, and \textbf{Foresight Guidance}~\citep{wang2025foresight} frames CFG as a fixed-point iteration with short inner loops. Complementary to these, \textbf{Prompt-Aware Guidance}~\citep{zhang2025prompt} adapts strength based on prompt complexity, \textbf{Learning-to-Guide}~\citep{galashov2025learn} employs meta-learning for optimal strategies, and \textbf{Saddle-Free Guidance}~\citep{yeats2025saddle} navigates the optimization landscape to avoid saddle points.
\textbf{Density Guidance}~\citep{zhao2025devil} extends flow matching by incorporating explicit log-density control to steer generation trajectories. While it offers a rigorous theoretical unification of prior heuristics via Score Alignment, the method is computationally expensive. Specifically, estimating the divergence requires Jacobian-Vector Products (JVP), which doubles the inference cost and introduces estimation variance.
% \textbf{Density Guidance}~\citep{zhao2025devil} introduces detail-aware generation by controlling the likelihood-detail trade-off. 
In Flow Matching models~\citep{esser2024scalingrectifiedflowtransformers, lipman2022flow, fan2025vchitect, liu2022flow, gao2024lumina}, \textbf{CFG-Zero*}~\citep{fan2025cfg0} compensates for velocity errors.
% and \textbf{FM-Auto-Driving} applies flow matching techniques to trajectory generation in autonomous driving.

\subsection{Challenges of Zero-SNR Sampling.} 
\citet{lin2024common} observe that guidance becomes unstable when sampling from true zero-SNR, attributing this to excessive update magnitudes. While Rectified Flow models~\citep{liu2022flow} resolve the training-inference mismatch of standard diffusion schedules by explicitly enforcing a pure Gaussian boundary, they remain susceptible to the zero-SNR guidance instability identified by \citet{lin2024common}. In this work, we analyze why early guidance update leads to instability. Based on this analysis, we propose an adaptive magnitude control mechanism to effectively stabilize the guidance trajectory.

\subsection{Challenges in Large-Scale Generative Models}

Early diffusion models, such as Stable Diffusion v2~\citep{rombach2022high}, operated on relatively compact latent spaces ($\approx 1.6 \times 10^4$ dimensions). In these settings, standard guidance strategies proved robust and forgiving to hyperparameter choices.

However, the transition to modern DiT-based architectures involves a massive increase in scale. For example, Flux~\citep{flux-2-2025} (generating 2K resolution) involves $\approx 10^6$ dimensions, and video models like Wan2.1 (14B)~\citep{wan2.1} exceed $10^7$ dimensions. \textbf{Empirically}, we observe that guidance strategies designed for smaller models become unstable at this scale. 
Building on the observation by \citet{lin2024common} regarding excessive guidance at Zero-SNR, we observe that this phenomenon is further exacerbated by model scale; specifically, the guidance update magnitude scales aggressively with the latent dimensionality.
% Specifically, the magnitude of the guidance update vector tends to scale aggressively with the model size~\citep{lin2024common}. 
Without careful regulation, these disproportionately large updates during the initial sampling steps lead to severe visual artifacts, such as color saturation and structural incoherence, effectively causing the model to ``overshoot'' the realistic image distribution. This necessitates a scalable, magnitude-aware correction mechanism.

